\scanpage{93}
\[
\text{b)}\qquad q(x+\underline i y)=0\qquad\text{i.e.}\quad q(x)-q(y)+\underline i\varphi(x,y)=0
\]
i.e.
\[
\left\{\begin{aligned}q(x)&=q(y)\\\varphi(x,y)&=0\end{aligned}\right.\tag*{(17)}
\]

Il faut faire agir là-dessus $G_m^{S[i]}$, qui s'identifie aux $a+ib$ tels que $\underbrace{a^2+b^2}_{\substack{\Vert\\\operatorname{Nm}(a+ib)}}\in\Gamma(S,\underline O_S^\times)$, \unsure{opérant|opèrent} sur $x+iy$ par
\[
\underbrace{(a+ib)}_{\alpha}\underbrace{(x+iy)}_{\unsure{\zeta|\xi}}=(ax-by)+i(bx+ay)
\]

Considérons l'ouvert \inserted{$\widetilde U$} de $\widetilde X^{S[i]}$ défini par
\[
\deleted{\unsure{[x,y]}}\quad q(x)=q(y)\in\Gamma(\deleted{\unsure{S}}-,G_m)\tag*{(18)}
\]
on a si $\zeta=x+\underline i y\in\Gamma\widetilde X^{S[i]}$ et $\alpha=a+ib\in\Gamma G_m^{S[i]}$
\[
\begin{aligned}q(\alpha\zeta)=q(ax-by)&=\unsure{q}a^2q(x)+b^2q(y)\unsure{-|+}2ab\underbrace{\varphi(x,y)}_{\substack{\Vert\\0}}\\&=(a^2+b^2)q(x)\in\Gamma G_m\end{aligned}
\]
donc $\widetilde U$ est stable par $G_m^{S[i]}$, et définit donc \deleted{\unsure{par}} par passage au quotient un ouvert :
\[
U\simeq\widetilde U/G_m^{S[i]}\subset X^{S[i]}\tag*{(19)}
\]

Considérons
\[
\deleted{\unsure{\varphi}}\qquad\widetilde U\longrightarrow\deleted{\unsure{Z}}\ \check{\mathbb V}(E)
\]
défini par
\[
\zeta=x+iy\ \overset{\widetilde f}{\longmapsto}\ \deleted{\unsure{[x,y]}}\ q(x)^{-1}[x,y]\tag*{(20)}
\]
% archive sheet number: 92
